\section{Introduction}
\label{sec:intro}

The chest radiograph is the most frequently acquired imaging examination and the initial test for pneumonia, pleural effusion, pulmonary nodules and tuberculosis (TB). Readers are scarce in many of the settings where TB and pneumonia are most common, which has driven sustained work on automated classification, from early TB screening on small curated sets~\cite{jaeger2014two,candemir2014lung} to deep networks trained on large hospital collections~\cite{wang2017chestxray8,rajpurkar2017chexnet,irvin2019chexpert,johnson2019mimic} and the many COVID-19 classifiers that followed~\cite{wang2020covidnet,apostolopoulos2020covid,chowdhury2020can}. No single public collection covers all classes of interest, so multi-class classifiers are usually trained on images pooled from several collections. Each collection differs in population, acquisition hardware, image processing and label provenance. Pooling therefore makes three failure modes likely: patient leakage between splits, proxy labels, and source confounding, in which a class comes mainly from one collection and the network can recognise the collection instead of the disease~\cite{zech2018variable,geirhos2020shortcut,degrave2021ai,maguolo2020critic}. Systematic reviews found these problems widespread among COVID-19 imaging models~\cite{wynants2020prediction,roberts2021common}, data leakage affects machine-learning-based science more broadly~\cite{kapoor2023leakage}, and models tend to perform worse on external than on internal data~\cite{yu2022external}. Reporting guidelines now ask for explicit descriptions of data sources, partitioning and external evaluation~\cite{mongan2020claim,collins2024tripodai}.

\paragraph{Related work} Source identifiability is well documented. Zech et al.\ showed that networks could recover the originating hospital system from a radiograph~\cite{zech2018variable}, Maguolo and Nanni that classifiers still separated COVID-19 images when the lung region was blacked out, which they read as dataset-specific patterns~\cite{maguolo2020critic}, and Dack and Dai that a ResNet-50 recovers the source dataset of chest radiographs with an F1 score near 99 and that this falls when pixels are shuffled~\cite{dack2025bias}. DeGrave et al.\ found that COVID-19 detectors relied on confounding factors rather than pathology and were likely to fail at new hospitals~\cite{degrave2021ai}, and Compton et al.\ showed that adding a second hospital can lower worst-group accuracy through hospital-specific artefacts~\cite{compton2023more}. At the level of the whole image, Zhang et al.\ identified contrast and sharpness as acquisition shortcuts in COVID-19 data~\cite{zhang2023certified}. Frequency-domain analyses are closer to the present question. Halimi et al.\ clustered the radially averaged spectra of chest radiographs and located failure modes of a classifier in particular spectral clusters~\cite{halimi2026spectral}, and Lu et al.\ related the spectra of model gradients to artificial frequency shortcuts in computed tomography and magnetic resonance imaging~\cite{lu2025spectral}. Kim et al.\ harmonised multi-institutional chest radiographs with a Laplacian-pyramid (multi-frequency) standardisation~\cite{kim2023homogenization}. None of the texts we read reports a per-band comparison of source and disease information in network features. Linear erasure is a published family of methods: iterative nullspace projection (INLP)~\cite{ravfogel2020null} and the closed-form LEACE~\cite{belrose2023leace}. Rahbar applied INLP to site information in frozen brain-MRI features and found that removal was destructive for dense segmentation~\cite{rahbar2026site}, and Bouaziz and Chebouba and Trivedi et al.\ reduced site or source information in chest radiograph features with adversarial or disentanglement training~\cite{bouaziz2026site,robinson2022shortcut}. Holding out one source at a time follows Yao et al.~\cite{yao2019strong}. Oakden-Rayner documented label problems in large public collections~\cite{oakdenrayner2020exploring}.

This study began as an audit of our own earlier pipeline (V5), which pooled eight public sources and recorded a high internal accuracy (\RVFiveAcc\% on $N=\RVFiveN$ test images). The audit found test images that shared a patient with training images, proxy or latent-disease labels, and external specificity of \RVFiveExtSpecK/\RVFiveExtSpecN. We rebuilt the data from official sources under rules fixed before training (V6) and evaluated a standard DenseNet-121 on it with external and per-source analyses. That rebuild gave a leakage-controlled benchmark on which a further question could be asked cleanly: where in the image does the source information that makes shortcuts possible reside, and what does removing it cost?

\paragraph{Contributions}
\begin{enumerate}
\item \textbf{A leakage-controlled, auditable six-class benchmark} from official sources, with per-file provenance, labelling rules fixed before training, patient- and duplicate-grouped splits, an external cohort, and an audit of the earlier pipeline (Sections~\ref{sec:materials} and~\ref{sec:disc-audit}).
\item \textbf{A band-by-band map of source and disease information (C1)} in frozen DenseNet-121 features of the cached images. To our knowledge, within the searches described below, we are not aware of prior work that quantifies, band by band, source versus disease information in CNN features of chest radiographs. The result is a null one: the source signal is present in every band and no single band carries it (Section~\ref{sec:res-c1}).
\item \textbf{A diagnostic application of published linear erasure (C3)} to the same features, showing how much source information INLP and LEACE remove and what that costs in disease information (Section~\ref{sec:res-c3}). We apply existing methods and do not claim a new one.
\item \textbf{A leave-one-source-out matrix (C4)} of linear heads on the benchmark, following Yao et al.\ and tied to the source probe (Section~\ref{sec:res-c4}). The protocol is not claimed as a contribution.
\item \textbf{A recorded negative decision (C2).} A band-stop retraining planned as a remedy was not run, because the go/no-go rule fixed in advance returned NO-GO (Section~\ref{sec:res-c2}). We regard the rule and its outcome as part of the design, not as a gap.
\end{enumerate}

The searches behind the novelty statement used web searches, Europe PMC, arXiv and Crossref in November 2026; they were not exhaustive and covered English-language sources only. Two of the closest works (Halimi et al.; Bouaziz and Chebouba) are preprints whose abstract pages were read, not the full texts. We propose no network architecture and make no claim of clinical validity. The classifier is a standard ImageNet-pretrained DenseNet-121~\cite{huang2017densely,deng2009imagenet}, so that data construction and evaluation are not confounded with architectural novelty. The main results come from one training run on a laptop CPU, with seed variability estimated from three runs of the selected build.

Section~\ref{sec:materials} describes the data, Section~\ref{sec:methods} the model and the C1--C4 methods with their decision rules, Section~\ref{sec:design} the experimental design, Section~\ref{sec:results} the results, Section~\ref{sec:discussion} the discussion and limitations, and Section~\ref{sec:conclusion} concludes. Tables moved out of the main text are in the supplementary file.
